.mt9
.mb7
.po8
.ll65
.hm4
.bp
.ig
CHAPTER 1
H O W	T O   U S E   T H I S	M A N U A L
.pp5
For best results read through this manual thoroughly at least once before
operating the system and do the practice run outlined in Chapter 4.
.pp
The Inventory Control System documentation includes two manuals.  This
manual is a Users' Manual which provides all the background information
and instructions necessary for successful operation.  The other is a
Specifications Manual which contains technical information on file accessing
and file descriptions.	The Specifications Manual is available only
with the purchase of the system itself.
.pp
The Users' Manual divides into an Introduction, a Reference Section, an
Operations Section, and an Appendix.  The Reference sections deal with
general concepts and procedures.  The Operations sections give
additional specifics on running the programs.
.pp
While Structured Systems Group is glad to provide helpful support to
registered owners of its products, if you have a problem, please make certain
you have read and understood the documentation and checked for hardware
malfunctions before giving us a call.
.he
.bp
.ig
CHAPTER 2
F E A T U R E S   A N D   F A C I L I T I E S
.pp5
The Structured Systems Group Inventory Control System provides small
to medium businesses with up-to-date and accurate information on the
quantity, value, and activity of their inventory items.  The Inventory
Control System has been designed to serve as the central program module
of a complete inventory management system.
The present
Inventory Control System is the first component of that complete inventory
package.  It
stands able to meet the full needs of many
businesses.  The system performs all the basic inventory control functions
such as adding to and depleting stock, adding new items to stock,
deleting items no longer carried, keeping track of quantities on hand,
quantities on order, quantities back-ordered, and inventory values.  Access
to file records by item number is quick, allowing timely management information
and fast customer service.
.he Introduction: Chapter 2
.pp
A special "auditability option" lets you create an ongoing hardcopy
record of stock additions, depletions, or both
additions and depletions.
If data is accidentally entered incorrectly, or if discrepancies
arise between your accounting and inventory records, the hardcopy transaction
proofs simplify the job of tracking down the errors.
.pp
The Inventory Control System generates a number of informative reports,
including:
.sp
.in4
.ti-4
1.  A detailed printout of the master item file, called the ITEM LIST,
that includes information on vendor number, retail price, average unit
value, most recent unit cost, quantity in stock, reorder point, reorder
quantity, date ordered, date due, quantity on order, and quantity back-ordered.
.sp
.ti-4
2.  A STOCK ACTIVITY REPORT that details the current period or period to date
activity for all or a selected range of stock items.  In addition to current
quantity on hand, retail price, unit cost, and stock value by item, the
report displays quantities added, quantities depleted, and the net change in
stock.
.sp
.ti-4
3.  A STOCK VALUATION REPORT that gives the average value of the item, the
replacement cost,
stockroom quantity, stockroom value, retail price for each item, and the
total inventory value for all items included in the report.
.bp
.ti-4
4.  A REORDER REPORT that automatically selects all out-of-stock items,
items below the reorder point, or items on order, as well as a range of
these items.
.qi
.sp
.ad
See Chapter 15 for sample reports.
.pp
The Inventory Control System incorporates SSG's Display Management System
(DMS), a full-screen formatting utility.  Like all Structured Systems Group
packages the Inventory Control System employs the most advanced human
engineering techniques.  Defensive programming ensures successful system
operation and trauma-free learning.  Edited input and logical commands
make for the clearest and easiest operation possible.
.pp
The Inventory Control System permits up to four disk drives and three item files.
The maximum number of inventory item records supported by the system is 32,767.
.he
.bp
.ig
CHAPTER 3
O V E R V I E W   O F	S E T U P   A N D   O P E R A T I O N
.pp5
Because the Inventory Control System is menu driven, you run the programs
by selecting them from the system menu.  The Inventory Control System also
employs SSG's Display Management System, a screen formatting system that
simplifies data and parameter entry. By moving the CRT cursor from field
to field within the display you enter or change the values in the fields.
The system edits your input, and issues descriptive error messages to ensure
smooth operation.
.he Introduction: Chapter 3
.sp
.ce
SYSTEM SETUP
.pp
The first step in system setup is to make a "system master" and "work disk"
from the "distribution disk" received from SSG.  The second step in system
setup is to define the characteristics of your CRT to the Display
Management System.  Step three of system setup is to set the system parameters.
There are four sets of parameters:
.sp
.in8
.ti-4
1.  General System Parameters:	These determine the company name, page
length, system date, date format (American or international), and whether
to sound the console "bell" if an error occurs.
.sp
.ti-4
2.  Item File Parameters:  These establish the number of item files (up to
3 are allowed), their drive locations, and the starting item number for
each item file.
.sp
.ti-4
3.  Audit File Parameters:  These determine whether the auditability option
is to be used, and if chosen, establishes whether the audit file is for
recording additions or depletions, or both, and which drive location holds
the audit file.
.sp
.ti-4
4.  Sort File Parameters:  These establish the input, workfile, and output
drive locations for sorting the item files.
.qi
.pp
Finally, step four is to actually build the item file by entering information
for each inventory item.
.bp
.sp
.ce
NORMAL OPERATION
.pp
Once your system is set up and the item file built, normal operation involves
adding quantities received, depleting quantities moved out of stock, and
producing the reports.
.pp
You can access inventory records from the item list in three ways:
.sp
.in8
.ti-4
1.  By entering the drive name and record number relative to the beginning
of the file on that drive (i.e., the 4th record on drive B).
.sp
.ti-4
2.  By entering the item number.
.sp
.ti-4
3.  Sequentially, by a next-record, or previous-record command.
.qi
.pp
Besides entering the quantity to be added to stock, you can also change
the unit cost, date ordered, date expected, quantity on order, and quantity
back-ordered.  If you want it to, the program automatically recomputes the
inventory value for that item by the weighted average method.
.pp
To deplete stock you enter the quantity leaving your stock.  If you desire,
the program
automatically recomputes the value of the remaining inventory by
subtracting the value of the items leaving stock.  This value is determined
by multiplying the average value of a stock item by the number of units
leaving stock.
.pp
If you choose not to calculate the value of your inventory by the weighted
average method, yet still wish to keep track of your inventory value,
you must calculate and change the inventory value manually.
.pp
Should you mistakenly add more units to stock than were actually received,
you can fix the error by reversing the transaction.  Reversing a transaction
is handled by adding a negative amount to stock, not by depleting the stock
in the amount of the error.  This method of "backing out" a transaction
maintains the accuracy of the stock activity report.  Erroneous depletions
are handled similarly.
.pp
If you choose the auditability option, normal operation also involves
generating TRANSACTION PROOF REPORTS after each session of adding to stock
or depleting stock.  The transaction proof reports are a hardcopy record
of the actual entries to the system, and thus serve as a check on the
accuracy of your entries, or a way to trace bookkeeping discrepancies.
Each transaction proof is numbered and dated to provide a continuous record of
your entries to the system.
.pp
At the end of each current or to-date period, the START A NEW ACTIVITY PERIOD
program lets you zero out the activity registers.
.sp2
.ce
REPORT PRODUCTION
.pp
Generating reports is a matter of choosing the report from the system
menu and telling the system which report options you want to use (such
as reporting only on a range of inventory items for example).  Chapter
15 presents examples of the reports available from the Inventory Control
System.
.sp2
.ce
FILE MAINTENANCE
.pp
When a new item is added to the item file, sorting may be needed,
depending on the item number.  After a sort, the program requires you
to check for duplicate item numbers by running the ITEM FILE
PRINT program which
can be instructed to print the entire item list with duplicate items flagged,
or to print only duplicates if found.
.pp
Duplicate file items or items no longer in stock can be removed by a delete
command.  After deletions are made, the ITEM FILE RECOPY program physically
removes the deleted records from the file,thereby compacting it to save disk
space.
.pp
You can change the item file parameters, for example, to inform the system
that additional drives are to be used, or to redistribute the inventory
items onto different drives.
.he
.bp
.sp8
.ce
[ GRAPHIC ]
.bp
.ig
CHAPTER 4
P R A C T I C E    R U N
.pp5
This Practice Run section provides a quick hands-on introduction to the
Inventory Control System.  Follow the instructions carefully and
make the entries exactly
as directed.  The Practice Run covers setting up a system
master (or program) disk, implementing the Display Management System,
establishing system parameters, building an item file, adding to stock,
depleting stock, and generating reports.
.he Introduction: Chapter 4
.pp
First, create a "system master" from the "distribution disk" you received
from SSG.  Place a formatted disk containing the CP/M operating
system, the CBASIC runtime interpreter program CRUN2, and the CP/M
PIP program on drive A (see the CP/M Features and Facilities manual
for how to SYSGEN an operating system onto the disk).  Place the SSG
distribution disk (write-protected)
on drive B and hit the RESET switch on your machine.
When the CP/M prompt character is showing (A>), type:
.sp
.ti10
PIP A:=B:*.*[OV]
.sp
.ad
and follow with RETURN; this command copies the programs on drive B onto
the disk on drive A.  When the PIP ends, remove the distribution disk
from B and store it in a safe place.
.pp
Now you will set up the Display Management System for your CRT. The
Inventory Control System comes ready to run on a Hazeltine 1500. If you
have that device, you need only erase the DEF type files from your system
master (or "program") disk (not from the SSG distribution disk!). Type
the command,
.sp
.in10
ERA *.DEF
.qi
.sp
.ad
and continue.
.pp
If you have one of the CRT's
listed in Chapter 17, locate the file that corresponds to your CRT and
then rename that file to a file called CRT.DEF by the command,
.sp
.ti10
REN A:CRT.DEF=A:nnnnnnnn.ttt
.sp
.ad
where nnnnnnnn.ttt is the filename and filetype of your CRT definition file.
You may then erase the remaining DEF files by the CP/M ERA command (be
sure not to erase CRT.DEF).  If you have neither a Hazeltine 1500 nor one of
the CRT's listed in Chapter 17, see Chapter 28 for instructions on how
to run the CRTDEF program to define your CRT to the DMS.
.pp
You should now have a system master with the Inventory Control System
programs on
drive A.  Insert an empty, formatted disk for a data disk on drive B.
Give the CTRL-C command or hit the RESET switch on your machine to indicate
that a new disk has been inserted.  You are now ready to begin operating
the system.
.pp
To begin, enter:
.sp
.ti10
CRUN2 IY
.sp
.ad
at the CP/M prompt character, and follow with RETURN.
.pp
Before continuing, the system creates a default parameter file.  When it
asks you to ENTER SYSTEM DRIVE, enter A followed by RETURN; this is the
drive where all parameter files will be located.
.pp
First up is the PARAMETER ENTRY MENU (see Chapter 27 for reproductions
of the screen formats).  At the SELECT A SCREEN request enter 1 followed
by RETURN to bring up the GENERAL PARAMETERS screen.  Make no changes
to the values shown on this screen, except to enter your company name, if
desired.  The CRT cursor should already be at the company name field.
.pp
If you make an error while keying in any entry to the system the command
CTRL-H backs the cursor up so you can correct the entry.  Enter ESCAPE
at any point to go from the General Parameters screen to the Parameter
Entry menu; enter ESCAPE once more to continue.
.pp
When asked for the date enter 2/1/79.  Follow with
RETURN.  The ITEM FILE PARAMETERS screen which appears next displays the
defaults you will use.	Make no changes or entries.  Enter ESCAPE to go
on to the main system menu.
.pp
At the SYSTEM MENU select program 9, CREATE SORT PARAMETER FILES, and
follow with RETURN.  After a message tells which drive holds the file
to be sorted (the input drive), give the temporary workfile drive as
drive B and follow with RETURN.  Enter N (for NO) followed by RETURN
when asked if you want to switch disks.  When asked if you want to PRINT
SORT PARAMETERS? answer YES, then RETURN.
Turn on your printer, align the paper, and
hit RETURN to print.
.pp
When the system menu returns, select program 10, ITEM ENTRY, and hit RETURN.
When the cursor is in the Item Number field, enter P-1, for item number P-1.
Hit RETURN.  After the program displays the ACCESS KEY, the cursor
moves to the description field.  Make entries for the item by filling in
the blanks with the information in the table below, hitting RETURN to advance
to the next field if a field is not filled. (Use CTRL-B to back the cursor
up one field if necessary.)  After making all the entries for the item, hit
ESCAPE to begin making entries for the second item.  Enter P-2 in the item
number field and continue as before.  After you've entered information
for the last item, hit ESCAPE to return the cursor to the item number field.
Enter ESCAPE again to move the cursor to the access key field.	Enter
ESCAPE from there to return to the system menu.
.sp
.po5
.ll75
.cp9
.li
ITEM DESC.  STOCK REO REO RETAIL INV.  VEND. COST/ QTY DATE   DATE   B-
 #	    ROOM  PT. QTY  PRICE VALUE	 #   UNIT  ORD ORD    DUE    ORD
------------------------------------------------------------------------
P-1  hammer  12    6   18   9.55 60.00 00001  5.00
P-2  saw      0    5   12  18.99  0.00 00001 10.00  12 1/1/79 3/1/79  0
P-4  nails   20   30  100   1.50 20.00 00002  1.00 100 1/7/79 3/7/79  0
P-3  tacks   15   10   50   1.25  7.50 00002   .50
P-5  tape    12   15   25   1.75  8.00 00015   .75  25 1/7/79 3/7/79  0
P-6  glue    25   10   35   2.00 25.00 00015  1.00
.po8
.ll65
.pp
At the system menu, select program 11, ITEM FILE SORT.
When the sort ends, the system menu reappears.
Select program 12, ITEM FILE PRINT.  Answer YES to the question,
DO YOU WISH TO PRINT ITEM FILES?  and follow with RETURN (be sure
your printer is ready).
.pp
You've now completed the system setup and are ready to begin normal processing.
Choose program 1, ADD TO STOCK, from the system menu.  When the
ADDITIONS TO STOCK screen appears with the cursor in the item number
field (see Chapter 27 for a reproduction of the display format), enter
the dummy transaction data in the table below.
After entering the data for a
transaction hit ESCAPE to begin the next.  After entering the last addition
to stock, hit ESCAPE, and ESCAPE again to return to the system menu.
.cp6
.sp
.li
    ITEM    QTY      UNIT    QTY   DATE     DATE    QTY
     #	   ADDED     COST    ORD   ORD	    DUE    B-ORD
    -----------------------------------------------------
    P-2     12	    10.00	   NONE     NONE
    P-4     75	     1.00			     25
    P-5     25	     1.00	   NONE     NONE
.sp
.ad
Ready your printer, then select program 3, TRANSACTION PROOF, to print
the additions audit file, a record of your entries to the system.
.pp
When the system menu returns, select program 2, DEPLETE STOCK.	A
message that the existing audit file has been proofed, but is not of
the depletions type appears.  It asks if you wish to erase it.	Answer
YES, then RETURN.  When the DEPLETIONS FROM STOCK screen appears, enter
the data in the table below, hitting the RETURN key after entering the quantity
depleted to begin the next transaction.
After entering
the last transaction, hit ESCAPE to return to the system menu.
.sp
.cp8
.li
	ITEM	       QUANTITY
	 #	       DEPLETED
     -----------------------------
	P-1		  2
	P-3		  6
	P-4		 30
	P-5		 10
	P-6		 16
.pp
Now run program 3, TRANSACTION PROOF, to print the depletions audit file.
The system menu returns after the printout.
.pp
Print the REORDER REPORT by selecting program 4.  When asked to SELECT A
REPORT OPTION, key in 3 to indicate you want a report on items that
are understocked (whose stockroom quantity is lower than the reorder point).
First make sure your printer is turned on, with the paper adjusted.  The
system menu reappears after the printout.
.pp
Print an ACTIVITY REPORT by choosing program 6 from the system menu.
When asked to select the ACTIVITY REPORT OPTION, enter 1 to generate
a report that covers all items.  The system menu returns after the
report prints out.  Enter ESCAPE from there to end the Invetory Control
System.  That's all there is to it.
.he
.bp
.ig
CHAPTER 5
A D A P T I N G    O F F I C E	  P R O C E D U R E
.pp5
Adapting your office routine to accomodate computerized inventory control
is a minor process in most cases.  As with any endeavor, careful planning
and forethought smooths the process considerably.
This section suggests procedures for setting up, normal operation, using
the reports, and backing up for safety.
.he Introduction: Chapter 5
.sp
.ce
SETUP
.pp
From your company records or a physical inventory compile a list of the
items
you stock.  Arrange the items into logical groupings if possible.
Besides the normal operational advantages, such an arrangement means
the ACTIVITY, REORDER and VALUATION REPORTS cover related items when
you exercise the report range option.
.pp
Assign an item number to each item,
leaving numbers open between logical groupings
to allow for adding more items at a later time.
You may use your current numbering system,
or renumber the list to conform to the system's 10 character item number
limit.	A cross-reference between old and new numbers could be maintained,
either manually, or automatically with SSG's Analyst system.
.pp
Item numbers may contain the characters A through Z, the numbers 0 through
9, slashes (/), and dashes (-).  Be sure the slashes and dashes line up.
For example, the dashes in the following item numbers are aligned properly:
.sp
.li
	    A180-005-6
	    A180-005-7
	    A180-016-5
.sp
.ad
To enter the number A180-5-6 instead of A180-005-6 would be inproper.
In the examples above, the first character of the item number (A) gives
the warehouse location.  The next character (1) the department number,
the next two (80) give the dissection, the three characters following the
dash (005 and 016) are a product code, and the final number is a size code.
.pp
After compiling the complete list of inventory items and assigning
them numbers, make an accurate determination of the quantity on hand
and value for each, either by physical inventory or company records.
.pp
In addition, you should have retail price, cost per unit, reorder point,
and standard reorder quantity information for each item.  For items
currently on order you should also have information on the date ordered,
date due, quantity ordered, and, if applicable, quantity back-ordered.
.pp
You may assign a vendor number for each item.  The system is set up to
use the 5 digit vendor numbers of the SSG Accounts Payable system.  If
you buy the same item through several vendors, you can establish a
separate record for each vendor, or leave the vendor number blank.
.sp2
.ce
NORMAL OPERATION
.pp
Under normal circumstances, additions to stock can be keyed directly into
the system when they arrive, or batched for later processing.  You can
take the information  from packing slips, invoices, and other original
media, or prepare a form of your own. The packing slips and invoices
can then be routed through your operation normally.
If you choose the auditability option, the batch processing method (saving
up additions or depletions
for entry all at once) helps ensure the accuracy of your transaction
proof reports (see Chapter 13).
.pp
You can handle depletions from stock similarly.  Enter the depletions
when they occur, or batch them for later entry.  You can take the information
directly from customer orders, pull-sheets, or forms designed especially
for the task.
.pp
Reorder or backorder information can be entered directly from your purchase
order books, buy sheets, or other media.
.sp2
.ce
USING THE REPORTS
.pp
The reports generated by the Inventory Control System provide
management, accounting, and stockroom personnel with up-to-date
and accurate
information on availability, turnover, order status, inventory value,
as well as retail and replacement costs.
.pp
If you choose the system's auditability option,
after adding to stock or depleting it you must generate a TRANSACTION
PROOF as a check on your entries.
Since the TRANSACTION PROOF details either additions or depletions,
but not both on the same report, batch processing is recommended if
auditability is selected.
Batching the entries means you generate one report for
additions, and one report for depletions.
.pp
The Inventory Control System prepares the reports listed below at any time:
.sp
.in10
.ti-3
1. ITEM LIST: a multi-purpose listing of all inventory items for use by
management, buyers, counter, and
stockroom personnel.
.sp
.ti-3
2. ACTIVITY REPORT: most useful for the company buyers and marketing
staff.
.sp
.ti-3
3. STOCK VALUATION REPORT: useful information for the company accountant,
controller, or other management. In some states, a real time saver at tax
time.
.sp
.ti-3
4. REORDER REPORT: a real help to your buyer, and up-to-date information
on availability for counter personnel or salespeople.
.qi
.sp2
.ce
BACKING UP FOR SAFETY
.pp
In order to achieve greater storage, the Inventory Control System
cannot automatically back up item files.  Therefore, it is important
for you to have at least one up-to-date copy of the file at all
times.	The simplest procedure is to make exact copies of the disks
containing the item (or even parameter) files with the full-disk copy
program that should have been provided with your disk hardware.  If this
program is not available, use the PIP command, or the SDCOPY program
provided on the SSG distribution disk.
.pp
SDCOPY works only for copying single density 8" IBM 3740 standard
compatible disks on a CP/M system.  To run the program, type SDCOPY at
the CP/M prompt, and follow with RETURN.  When the message PUT SOURCE ON
A, DESTINATION ON B, THEN TYPE RETURN appears, you remove any disks
currently in the A and B drives and put the disk to be copied on A.  Put
an empty, formatted disk (or any disk you don't mind writing over) on B.
Hit RETURN to begin.  When the copy is complete the reboot message
PLACE SYSTEM DISK IN DRIVE A, TYPE RETURN TO REBOOT appears.  Remove the
two disks and insert the disk that was originally on A.
You can use the same disks over and over as the backup disks.
.pp
A printout of the ITEM LIST also provides a margin of safety.  Should your
hardware be out of operation for an extended period, or your files
destroyed, this hardcopy backup enables you to track the changes in
stock manually until the problem is taken care of.
.he
.bp
.ig
CHAPTER 6
O P E R A T I O N A L	 B A C K G R O U N D
.sp2
.ce
HARDWARE AND SOFTWARE REQUIREMENTS
.he Reference: Chapter 6
.sp
.pp5
The SSG Inventory Control System
will run on any 8080 or Z-80 based microcomputer
under the control of the CP/M operating system or a CP/M compatible
operating system.  Your microcomputer system must have the following
minimum hardware and software components:
.pp
1.  48K (kilobytes) of Random Access Memory (RAM).  (CP/M configured for
45K or larger.)
.pp
2.  Two floppy disk drives with a minimum capacity of 200K.
.pp
3.  A 132 column wide printer capable of skipping to the top of the next
form in response to the ASCII form-feed character.
.pp
4.  For the console device, a CRT (Cathode Ray Tube) with a direct cursor
addressing facility, a clear screen control, and
minimum dimensions
of 24 lines by 80 columns.
.pp
5.  System software including an operating system and the CBASIC2 language.
.pp
Many microcomputers can be expanded or adapted to meet these requirements.
Your local dealer can give you specific information about the hardware
currently available.
.sp3
.ce
THE CBASIC PROGRAMMING LANGUAGE
.sp
.pp
The programs which make up the Inventory Control System are written
in a dialect of the BASIC programming language, called CBASIC-2,
which is a product of Software Systems. (Questions relating to
the CBASIC language should be refered to Software Systems.)
.pp
Structured Systems Group distributes its programs in an intermediate
file format (with the filetype INT), which must be further translated
at run time by a
facility of the CBASIC language called the "interpreter".  The interpreter's
program name CRUN2 must precede any INT program for it to run.
For example, to run the Inventory Control System you would enter CRUN2 IY.
This run time interpreter comes as a part of the CBASIC
language package.
With the CBASIC language you
can also write programs yourself to extend the capabilities
of the Inventory Control
System, or to perform other data processing tasks.
.sp3
.ce
THE CP/M OPERATING SYSTEM
.sp
.pp
An operating system is a set of instructions which enables the computer to,
among other things, read and write information on floppy disks, move and
maintain files, and provide for console communication.	The Inventory Control
System
employs the CP/M operating system, or one compatible with CP/M.  CP/M (TM) is a
product of the Digital Research Company, and any questions concerning the
CP/M system specifically should be referred directly to that company.
.pp
CP/M comes documented by a series of six manuals.  The first of the six,
"Features and Facilities", contains useful and important information
relating to your operation of the Inventory Control System.  Be particularly
careful to read and understand the DIR, ERA, REN, PIP, SYSGEN, TYPE,
and STAT commands before using the Inventory Control System.
.sp3
.ce
LOGGING ON
.sp
.pp
When you first turn your computer on, its random access memory (RAM) is
completely empty.  In this state the computer can do nothing.
Usually, however,
a feature has literally been built into the microcomputer hardware which
directs it to look on a predetermined section of a certain disk drive to find
what to do next.  If everything is working normally, your computer should
look on the outer two tracks of the A disk, find the CP/M operating system
(if it has been SYSGEN'ed onto the disk; see the CP/M "Features and Facilities"
manual for details on how this is done), load it into the RAM, and then
display "A>" on the CRT, indicating it is ready to go to work.  This
process is called "booting" since the computer is, in effect, lifting itself
up by its own bootstraps.
.pp
The "A>" character is called a "prompt".  When the A> prompt appears on the
screen, the operating system is said to be "logged on" to the A disk drive.
While logged on to the A drive, you can call up programs or data files on that
disk by typing the program name.
To request a file on another disk (unless this is done automatically
by the program you are using) first log on to that disk by keying in
the drive name (the letter B on a two disk system), followed by a
colon (:) and a carriage return and then type the program name.
(The carriage return is simply the RETURN
key as it is on a typewriter; it is referred to as CR.)
.sp
.li
	For example,
.sp
.li
	   A>		   [system prompt]
	   A>B: 	   [user enters B:, then RETURN]
	   B>
.sp
.ad
To return to the A drive, enter:
.sp
.li
	   B>A: cr
.pp
Whenever the prompt appears you are operating under the CP/M
environment.  It goes away when you enter the Inventory Control System, and
returns when you leave it.
.sp3
.ce
FULL-DISK BACKUP
.sp
.pp
Caution and common sense dictate that whenever you purchase or receive a
new program or system  on a floppy disk, you immediately copy it
onto a fresh disk and store the original in a fire-proof, moisture-proof  box
to protect it against inevitable human error and unpredicable system failure.
With the original SSG distribution disk tucked safely away from the ravages of
dirt, grease, magnet, and mangle, you avoid the cost and inconvenience of
unnecessary replacement.  This warning applies doubly for disks
containing data crucial
to your operation.  Do not make the mistake of being "penny wise and pound
foolish"; the disk itself is almost always worth far less than the information
residing on it.
.pp
The manufacturer of your disk drive(s) should have provided a full-disk copy
program with the disk hardware.  If this program is not available, use the
PIP command that comes with the CP/M operating system, or the SDCOPY program
provided with the Inventory Control System and explained in the previous
chapter.
.pp
To use	PIP
place a fresh disk containing an operating system and the PIP program on
drive A, and the Inventory Control System
distribution disk on drive B.  Hit the reset
switch on your machine and then type
.sp
.li
	A>PIP A:=B:*.*[OV]
.sp
.ad
This will copy all the programs and files on B onto the A disk.
See the CP/M documentation supplied with
your system for a complete description of the PIP program.
.sp3
.ce
CHANGING DISKS BETWEEN PROGRAMS
.sp
.pp
Whenever  you change disks between programs it is of the utmost importance
that you enter a CTRL-C character immediately after inserting the new disk.
The
CTRL-C character, which appears on your screen as ^C, informs the computer
that a new disk with different data has been inserted.	This "bootstrap"
operation (or "re-booting") causes the CP/M system to reset its internal
maps of disk usage, and prevents it from writing over (thus destroying)
files which exist on the new disk.
.pp
To enter CTRL-C, or any CONTROL character, hold down the CONTROL key (CTRL
or CTL on some machines), while simultaneously pressing the letter C
(either upper or lower case).
Never key in a CTRL-C character unless the CP/M prompt is showing.  Doing
so may end the Inventory Control System, with
a probable loss of data.
.he
.bp
.ig
CHAPTER 7
SETTING UP THE DISPLAY MANAGEMENT SYSTEM
.pp5
The Structured Systems Group Display Management System is a full-screen
information display and data entry utility.
Chapter 27 of this manual shows each of the screen formats used by
the Inventory Control System.  Data and parameter entry  to the Inventory
Control System takes place by "filling in the blanks" on the screen.
The Display Management System
makes it possible for related information and requests to
appear on the screen together.	This makes for easy record or parameter
review and fast, accurate, data entry.
.pp
Because most CRT devices differ, the Display Management System must
be informed of the characteristics of your particular device.  These
characteristics are held in a file of the type DEF that is read by DMS.
There
are three different methods for defining your CRT to the DMS, depending
on your situation:
.sp
.in10
.ti-3
1. If you have a Hazeltine 1500 device, you need do nothing besides
erase all the DEF type files on your system master or program disk
(but not the SSG distribution disk!), as well as the program CRTDEF.INT.
.sp
.ti-3
2. If your CRT appears on the list below you must carry out a simple
REN (rename) command
(see Chapter 17), and then erase all the DEF
files on your system disk, except the one called CRT.DEF.  You may also erase
CRTDEF.INT.
.qi
.sp
.li
		    ADM 3A
		    BEEHIVE 150
		    BEEHIVE 152
		    BEEHIVE 157
		    CROMEMCO 3100
		    DYNABYTE NAKED TERMINAL
		    SOROC IQ 120
		    IMSAI VIOC
.sp
.in10
.ti-3
3. If you have neither a Hazeltine 1500, nor one of the CRT's above,
you must run a program provided on your disk called CRTDEF.  CRTDEF creates
a file by an interactive question an answer process
that defines your CRT to the Display Management System.
Chapter 28 gives complete instructions.
.qi
.he
.bp
.ig
CHAPTER 8
G E N E R A L	A N D	A U D I T   P A R A M E T E R S
.pp5
Parameters can be defined as constant elements or
factors within which actions or
choices are constrained.  Because not all users have the same needs, you can
adapt the Inventory Control System to your requirements
by changing parameter values, rather than by expensive and complex
reprogramming.	The Inventory Control System parameters divide into four
general types:
.he Reference: Chapter 8
.sp
.in7
.br
1. General Paramters
.br
2. Audit Parameters
.br
3. Item File Parameters
.br
4. Sort Parameters
.qi
.sp
.ad
This chapter covers the first two.  Chapter 9 covers the last two.
.pp
The General Parameters control such things as the company name, date format
(whether American or international), page width, lines to print per page,
or whether the console bell sounds when an error occurs.  Other parameters
classified as general parameters determine what time periods (week, month,
quarter) are used for accumulating current period stock activity statistics,
and what time periods (month, quarter, year) for accumulating to-date
activity statistics.  The average valuation parameter lets you decide
whether you want the system to calculate the value of your stock by the
weighted average method and print a VALUATION REPORT.  Another
parameter sets an array size according to the amount of random access
memory you have; the larger the array, the faster inventory items can
be retrieved from the file.
.pp
The Audit Parameters let you determine whether you want to keep a hardcopy
record of your addition and depletion entries to the system as a way to
trace accounting or inventory errors, and check the accuracy of your
entries.  You can order a record of additions to stock, a record of
depletions from stock, or both.  You can also state on which disk
drive to place the audit file which accumulates the entry information.
.pp
When you implement the Inventory Control System for the first time, you
move directly to the parameter entry screens where you can examine
the default parameter values and
change them as needed.	Later, you set the general and audit
parameters by selecting
the PARAMETER ENTRY program from the system menu.
.he
.bp
.ig
CHAPTER 9
ITEM FILE AND SORT PARAMETERS
.pp5
By changing the Item File and Sort Parameters the Inventory Control
System can be set up to accomodate different numbers
of disk drives, and different numbers of stock items.  If the number
of items you stock is greater than the number that fits on a single
disk, you will need to break up the item file and spread it over
several disks by changing the Item File parameters.  And, since access
to inventory items requires sorted item files, the Sort Parameters must
be set to correspond to the item file arrangement.
.he Reference: Chapter 9
.pp
What are the item file parameters?  The system needs to know how many
item files you want.  Up to 3 files are possible, with one item
file only per disk.  For each item file you specify, the Inventory
Control System must know its disk drive location.  If more
than one item file is desired, the system will want to know the lowest
(first) item number on the second file, and the lowest (first)
on the third file.
.pp
At first time startup, the program asks for the Item File parameters
automatically, right after you've set the General and Audit parameters.
Later, if the number of items you stock grows, or if you wish to
redistribute the inventory items to different item files, you can
change the Item File parameters through the program on the main system
menu called ITEM FILE RECOPY.  After you change the parameters, the
program creates any new item files, and redistributes the inventory
items according to the new parameters.
.pp
To sort an item file, the system must know where to locate it (the input
drive), where to build the temporary workfile it needs for the sort (the
workfile drive), and where to put the sorted file when finished (the
output drive).	The sort parameter files contain this information.
.bp
.pp
If your item files are large, or the number of
disk drives available is limited, you may need to remove the input file
disk and replace it with another disk to receive the sorted
output.  The Inventory Control System provides for this possibility;
the Sort Parameter Files inform the system that you want this switching
to occur.  The next chapter describes disk drive allocation and
sort procedures in more detail.
.pp
If your files increase in size or number, you may need to alter
the sort parameters.  You can do this by choosing the
CREATE SORT PARAMETER FILES program from the system menu.
At the end of the CREATE SORT PARAMETER FILES function you can print out
the sort parameters for a hardcopy reference.
.he
.bp
.ig
CHAPTER 10
DISK DRIVE ALLOCATION AND SORTING CONSIDERATIONS
.pp5
The disk drive you assign each item file depends on how many disk drives
you have, how many item files you have, and how large your item files are.
The possible configurations range from one item file on a two drive system,
to three item files on a four drive system.
Chapter 19 suggests the best arrangement
for each possibility.
.he Reference: Chapter 10
.pp
Files other than item files (such as the audit and INT files)
may occupy a significant amount of disk space,
and thus should be assigned drive locations with this in mind.	The
program (INT) files take up the most room, and are always assumed to be
on the currently logged drive. The two parameter files (IYPR1 and IYPR2)
and the one to three
sort parameter files (type SRT) take up little room, and reside on the
drive you specify at the beginning of the program as the "system drive".
The system drive must be either drive A or drive B.
The audit file may be large or small depending on how often you proof and
erase the
current batch of transactions and begin a new one.  You specify the
drive location of the audit file at the AUDIT PARAMETERS screen that comes
up when you choose the PARAMETER ENTRY program from the system menu.
.pp
An item file must be sorted only when you add an item to it out of order
or cause it to become disordered by changing an item number.
Only the file affected must be sorted.	For example, if item file
#1 begins with item number K400 and runs to item number M800, adding an
item with a number that falls in between (if, of course, such a number is
available; otherwise a duplication would occur) such as L200
necessitates a sort.
After the sort, you must check for duplicates by running the ITEM FILE
PRINT program.	Only the ITEM ENTRY, ITEM FILE PRINT, and ITEM FILE SORT
programs can be run when unsorted files exist.
.pp
Only item file #1 may be sorted from the menu, and then only if
disk space is available (on any drive) to hold the temporary workfiles
needed for the sort.  There need not be space available for the
sorted output, since you can arrange to replace the input (unsorted)
disk with another disk to receive the sorted output.
The disk used to receive the sorted output may not contain an item file
of any kind.
.bp
.pp
To sort the other item files, or item file #1 if it is large, you end
the Inventory Control System and use a special sorting disk (which
contains the QSORT program and the appropriate sort parameter files),
to run the sort outside the system.
A printout of the sort parameters you specified
(input location, workfile location, and output location
for each partfile) which you can generate
at the end of the CREATE SORT
PARAMETER FILES program
helps you keep track of where each disk should
be.  Complete details and a step-by-step explanation are found in
Chapter 22.
.he
.bp
.ig
CHAPTER 11
B U I L D I N G   T H E   I T E M   F I L E
.pp5
The item file is a master list of all items in your stock.  Each
item takes up one record on the file.  Each record contains the item
number, a description of the item, its cost, its retail price, the value
of the quantity on hand,
reorder information, vendor number, and information on activity.
.he Reference: Chapter 11
.pp
Building the item file is the final step in system setup.
If no valid records exist on the file, as is the case when you first install
your Inventory Control System, the system goes immediately into the adding
mode when you choose the ITEM ENTRY program from the menu.  Once valid
items exist on the file, a control command is needed to add additional records
to the file.
.pp
To add items to the file, choose the ITEM ENTRY program from the system
menu.  When the item entry screen (see Chapter 27) appears,
enter the adding mode by the appropriate control command, or begin adding
if this is the first time setup.
Once in the item entry adding mode you enter the item's number.  After
the program displays the item's drive location and record number relative
to the beginning of the file on that drive, you enter a description of the item
(for example, "Sport shirts, lg.") and the current stockroom quantity.
Enter the reorder point, reorder quantity, vendor number, cost per unit
(usually the last price paid), the quantity on order, the date ordered,
the date due in, and the quantity back ordered (if any).  After making
these entries you can move the cursor from field to field indefinitely by
RETURN, or begin the next record by hitting ESCAPE.
.pp
The ITEM ENTRY program is primarily for building the item file, but
can also be used to examine, change, or delete items from the file.  A SAVE
command is also available; this command updates the disk directory
so the entries you make are not lost should the program end prematurely
due to a hardware failure (such as a pulled plug).
.pp
You can access records to examine or delete them by entering the
item number in the ACCESS KEY field, by entering the drive and number
of the item relative to the beginning of the file
on that drive, or by a next-record
or previous-record command.
.bp
.pp
You may also use the ITEM ENTRY program to update or change the reorder
information, the item description, and vendor number.  If, however, the
stockroom quantity is altered through this program, the stock activity
report will not reflect this change, and will thus be inaccurate.  If
auditability has been requested you should not alter the inventory value
through this program since it would not be reflected in the transaction
proofs.
.he
.bp
.ig
CHAPTER 12
A D D I N G   A N D   D E P L E T I N G    S T O C K
.pp5
Once your system is set up and the master item file built, additions
to and depletions from stock are handled by the ADD TO STOCK and
DEPLETE STOCK menu functions.  When you select the ADD TO STOCK or
DEPLETE STOCK programs from the system menu the ADDITIONS TO STOCK
or DEPLETIONS FROM STOCK screens appear with the CRT cursor in the
ITEM NUMBER field (see Chapter 27 for reproductions of the DMS screen
formats).  You may then call up the appropriate records by one of three
methods:
.he Reference: Chapter 12
.sp
.in10
.ti-4
1.  Enter the item number.
.sp
.ti-4
2.  Enter the drive location and record number relative to the beginning
of that drive's item file.
.sp
.ti-4
3.  Give a next-record or previous-record command.
.qi
.pp
When the record appears on the STOCK ADDITIONS screen you can enter the
number of units received and their cost, and update the order information.
The program automatically adds to the stockroom quantity for the item,
subtracting like amounts from the quantity ordered and backordered
figures (the quantity back-ordered is treated as a subset of the quantity
on order).  If you've chosen the average valuation feature the program also
increases the inventory value by the amount derived from multiplying the
number of units times the unit cost.  It also adds the number of units to the
current period and to-date period activity registers for the item which are
the basis of the ACTIVITY REPORT.  If you've chosen to audit additions to
stock, the transaction is recorded in the audit file.
.pp
When items leave your stock you use the DEPLETE STOCK program to enter the
change.  Access to records works the same as for addition to stock.
When the record appears, you enter the number of units leaving your stock.
The DEPLETE STOCK program decreases the stockroom quantity figure and adds
to the current period and to-date period activity registers.
.pp
If you've chosen average valuation the DEPLETE STOCK program also decreases
the inventory value for the item based on the weighted average value of the
item.  For example, if 5 units were added at $10 each, 5 units at $12 each,
and 10 units at $20 each the weighted average is figured as follows:
.cp8
.sp
.li
	       5 @ $10 = $50
	       5 @  12 =  60
	      10 @  20 = 200
	     ----	-----
	      20	 310

	   Weighted average = 310/20 = $15.50 per unit
.sp
.ad
Should a depletion from stock occur now, the inventory value for the
item would decrease at the rate of $15.50 per unit.  The weighted average
is recomputed each time new items are added to stock.
.pp
If you've chosen to audit depletions from stock the transaction is recorded
in the audit file.  If you've ordered both additions and depletions
audited you must print a TRANSACTION PROOF of the current audit file,
and then erase it, before you can go from the ADD TO STOCK program to
the DEPLETE STOCK program or vice versa.
The transaction proof is a hardcopy record of your addition and
depletion entries and their effect on inventory value.	You generate the
transaction proof by choosing the TRANSACTION PROOF program from the
system menu.   With the auditability option, batching transactions
(grouping them for entry all at once) is best, although you could move
back and forth from addtion to depletions as long as you printed the
transaction proof and erased the audit file
each time.  See Chapter 13 for more information on
the auditability option and transaction proofs.
.pp
Additions and Depletions should only be made through the addition and
depletion programs, not through the ITEM ENTRY program.  This keeps
the activity registers accurate, and records the transactions in
the audit file.
.pp
If the number of units added to stock is too large,
you must correct the mistake by making a
negative addition to stock in the amount
of the error, not by making an offsetting depletion.
Similarly, if the amount depleted is too large, you must "back out" of
the error by making a negative depletion, not an offsetting addition.
This makes sure the activity registers do not reflect additions and depletions
that did not actually take place.
.pp
If you choose to track the value of your inventory by the average
valuation facility, the amount subtracted from the inventory value (when
backing out an addition to stock error), or the amount added (when backing
out a depletion error) could differ from the original amount if the average
value for the item has changed.  Allowing the audit file to contain only
add to stock transactions or deplete stock transactions, but not both
at the same time, increases the chance that you will catch the error on the
on the transaction proof
before the average value of the item is altered by an intervening addition
at a different cost per unit.
.he
.bp
.ig
CHAPTER 13
A U D I T A B I L I T Y    A N D   T R A N S A C T I O N
P R O O F S
.pp5
The auditability option of the Inventory Control System allows the user to
leave a hardcopy inventory transaction history.  In addition to its use
in documenting sound financial and managerial practice, this hardcopy trail
facilitates error tracing and recovery.  It provides an ongoing check on
the accuracy of your inventory control system, making your automated
inventory and your physical inventory match as closely as possible.
.he Reference: Chapter 13
.pp5
The hardcopy printout is called a TRANSACTION PROOF.  The transaction proof
shows the item added or depleted, the quantity, the unit cost, and the
change in stock value brought about by the transaction.  It also shows
addition or depletion errors that were backed out as cancelled transactions.
.pp
You can order additions only, depletions only, or both additions and
depletions to be audited.  You make these choices when setting the
system's AUDIT PARAMETERS. A special audit file is created at the beginning
of the ADD TO STOCK or DEPLETE STOCK program, and resides
on the disk drive you specify in the AUDIT PARAMETERS until erased.
.pp
If you begin by adding to stock, the audit file created will be of the
additions type and will only accept add to stock transactions.	To maintain
the accuracy of your transaction history you must
print a transaction proof covering all additions transactions and then
erase the additions audit file before you choose the DEPLETE STOCK
program from the menu.	If you try to use the deplete stock program
while the additions audit file exists, a message informs you that the
current audit file is of the additions type. The message says whether all the
transactions in the file have been proofed, and asks if you wish to erase it.
.pp
If you begin by making depletions, a depletions audit file is created
and likewise must be proofed and erased before you return to
the add to stock program.
.bp
.pp
Of course, you can choose not to proof the audit file, erase it, and
go on to the desired function.	Doing so, however, means the transaction proofs
will not provide a complete and accurate record of your entries to the
system, and thus somewhat defeats the purpose of the auditability feature.
.pp
Each time an audit file is erased and a new one created the batch number
that identifies the audit file increases by one.  Each time an
audit file is proofed, the proof number for that batch increases by one.
Thus you can tell by looking at the audit proofs that a particular
audit file was never proofed since a batch number would be missing.
However, even if the batch numbers are contiguous the audit trail may not
be complete, since a batch may have been proofed at least once, but not
after additional transactions have been entered.  You should always
print a transaction proof before erasing a file.
.he
.bp
.ig
CHAPTER 14
A D D I N G ,  D E L E T I N G ,  A N D   C H A N G I N G
F I L E    I T E M S
.pp5
As you add new items to your stock and remove others no longer carried,
you'll want to add or delete item file records.  Your may also
want to alter the reorder point, reorder quantity, retail price, vendor
number, cost per unit, or update
order dates and quantities for an item.  These
additions, deletions, and changes are made through the ITEM ENTRY
program.
.pp
All records added to the file are placed at the end of the file regardless
of the item number of the record.  If the item number of the new inventory
item is greater than the item number of the last record on the file, the
file remains in sorted order.  If it is less, adding the record puts the file
out of order.
The program issues a message that if the item is
added the file will need to be sorted.	After sorting, to make sure
the file does not contain items with duplicate numbers, you must run
the ITEM FILE PRINT program.  This program prints out the entire item
list, marking any duplicates with asterisks, or, if you desire, prints
out only duplicate records that it finds.
.pp
To add a new item you call the ITEM ENTRY program from the menu and give
the command for adding new items.  Enter the item number and then
the information for the item.
.pp
When you first build the item file, you might want to leave certain
item numbers unused so new items of the same logical type will
be placed on the file at appropriate points.  Doing so will mean that
when you choose the option of reporting on a range of items, the range
will include related items.
.pp
To delete an item no longer carried in stock, you call the ITEM ENTRY program
from the menu, bring up the record, and issue the delete command.
Although the deleted record remains physically on the file until you
run the ITEM FILE RECOPY program, you cannot access or report on it.
.pp
To change an item on the file, call the ITEM ENTRY program from the
menu,  bring up the record and change the desired fields.  If you have
chosen the auditability option, never change
the stockroom quantity field through the ITEM ENTRY program; use the
ADD TO STOCK or DEPLETE STOCK programs.
If you've chosen to use the average
valuation feature, make no changes to the inventory value field; it
adjusts automatically whenever stock is added or depleted.  Changing
the item number may put the file out of order and necessitate a sort.
.he
.bp
.ig
CHAPTER 15
P R O D U C I N G    T H E    R E P O R T S
.pp5
You can order any of the four reports produced by the Inventory Control
System at any time by choosing them from the menu.  You can order
an ITEM LIST, a REORDER REPORT, a STOCK ACTIVITY REPORT, or a
stock VALUATION REPORT.
.he Reference: Chapter 15
.pp
To print the ITEM LIST choose the ITEM FILE PRINT program from the
system menu.  You can print the entire list, or only a selected
range of items.  The ITEM LIST shows the disk drive location, relative record
number, item number, item description, vendor number, retail price,
inventory value, cost per unit, number in stock, reorder point, reorder
quantity, date ordered, and quantity back-ordered.  In addition, this
report states the total number of records, the total number of deleted
records, and the drive location for each item file.
.pp
Using the REORDER REPORT you can report on items whose stockroom
quantity has fallen below the reorder point, on items out of stock
completely, or on items currently on order.  For each of these reorder report
options you can order a report on all, or only on a range of
items.
.pp
You can order an ACTIVITY REPORT on all items, on items active
in the current period (week, month, or quarter), or on items active
during the month, quarter, or year-to-date period.  For each activity
report option you can choose to report on all, or only a selected range
of items.  At the end of the current period (week, month, or quarter),
or at the end of the to-date period (month, quarter, or year)
you must reset the activity registers to zero.	To do this you run
the START AN ACTIVITY PERIOD program.
.pp
The inventory VALUATION REPORT is available whether or not you have chosen
to use the average valuation facility.	However, if you do not choose
average valuation, the report will be  meaningful only if you manually
compute the changes in inventory value and enter them through the ITEM ENTRY
program.
You can order the report on all file items, or only a selected range.
In addition to giving the inventory value by item, the report also
prints a total value for all items.
.pp
Representative report samples appear on the following pages.
#
